\chapter{Backend приложение и его основные функции}
\label{ch:backend}
\section{Лабораторная работа №9. Backend-сервер на C++}

\subsection{Теория: ZMQ и передача данных между Android и сервером}

ZeroMQ - это библиотека для асинхронного обмена сообщениями между процессами или по сети. Она работает поверх обычных сокетов, но упрощает программирование: не нужно управлять соединениями, повторными отправками и буферизацией. В проекте используется шаблон REQ/REP. Android выступает в роли клиента, отправляя JSON с координатами на сервер, а сервер отвечает коротким подтверждением. Данные передаются по TCP. Сервер привязан к порту 5555 на всех интерфейсах, клиент подключается по IP-адресу компьютера в той же локальной сети \cite{zeromq}.

\subsection{Реализация}

Сервер написан на C++ с использованием библиотек: libzmq для сокетов, nlohmann/json для парсинга JSON, ImGui и ImPlot для графического интерфейса, SDL2 как бэкенд окна.

\subsubsection{Структуры данных}

Для передачи данных из серверного потока в GUI-поток определена структура LocationData с координатами, высотой и временем, а также мьютекс для потокобезопасного доступа.

\begin{lstlisting}[language=C++, caption=Структура LocationData, label=lst:locationdata]
struct LocationData{
    float latitude=0.0f;
    float longitude=0.0f;
    float altitude=0.0f;
    float accuracy=0.0f;
    std::string time="No data";
    long long time_milliseconds=0;
    TrafficData traffic;
    std::vector<CellTowerData> cellTowers;
    std::mutex mutex;
};
\end{lstlisting}

\subsubsection{Серверный поток}

Функция run\_zmq\_server() создаёт ZMQ-контекст и сокет типа REP, привязывает его к порту 443. В бесконечном цикле сокет ожидает сообщение от клиента. При получении JSON-строки она парсится, извлекаются поля с координатами, временем и данными о сотах. Для потокобезопасного доступа используется мьютекс. Сервер также сохраняет данные в JSON-файл и базу данных.

\begin{lstlisting}[language=C++, caption=Обработка входящих данных на сервере, label=lst:server_parse]
void run_zmq_server(LocationData* loc){
    zmq::context_t context(1);
    zmq::socket_t socket(context,ZMQ_REP);
    socket.bind("tcp://*:443");

    while(true){
        zmq::message_t request;
        socket.recv(request);
        std::string msg=static_cast<char*>(request.data());

        try{
            json j=json::parse(msg);
            {   std::lock_guard<std::mutex> lock(loc->mutex);
                if (j.contains("location")&&j["location"].is_object()){
                    auto& loc_data=j["location"];
                    loc->latitude=getFloatValue(loc_data,"latitude",0.0f);
                    loc->longitude=getFloatValue(loc_data,"longitude",0.0f);
                    loc->altitude=getFloatValue(loc_data,"altitude",0.0f);
                    loc->accuracy=getFloatValue(loc_data,"accuracy",0.0f);
                    loc->time=loc_data.value("time","No data");}}
            g_measurementCounter++;
            saveToJsonFile(*loc,g_measurementCounter);
            updateSignalHistory(*loc);}
        catch(const json::parse_error& e){std::cerr<<"Ошибка JSON: "<<e.what()<<std::endl;}}}
\end{lstlisting}

\subsubsection{GUI-поток}

Функция run\_gui() инициализирует окно SDL2, контекст OpenGL, ImGui и ImPlot. В каждом кадре отображается окно с заголовком LOCATION \& TELEPHONY, где выводятся значения из структуры LocationData.

\subsubsection{Запуск потоков}

В функции main() создаются два потока: gui\_thread и server\_thread. Им передаётся указатель на общую структуру LocationData.

\begin{lstlisting}[language=C++, caption=Запуск потоков в main(), label=lst:main_threads]
int main(){
    if (g_database.connect()){std::cout<<"\033[32mПодключение к БД успешно\033[0m"<<std::endl;}
    std::thread gui_thread(run_gui,&g_locationData);
    std::thread server_thread(run_zmq_server,&g_locationData);
    gui_thread.join();
    server_thread.join();}
\end{lstlisting}

\begin{figure}[h]
\centering
\includegraphics[width=0.7\textwidth]{lab9_backend.png}
\caption{Окно сервера на C++ 1 версии}
\end{figure}

\section{Лабораторная работа №13. Работа с базой данных PostgreSQL}

\subsection{Теория: база данных и её схема}

PostgreSQL - реляционная система управления базами данных. В проекте используется для хранения всех измерений, поступающих от Android-приложения. Схема базы данных состоит из двух таблиц: based (основные измерения) и cells (данные о сотах). Связь один-ко-многим: одно измерение может содержать информацию о нескольких вышках сотовой связи.

\begin{lstlisting}[language=SQL, caption=Схема базы данных PostgreSQL, label=lst:db_schema]
CREATE TABLE based (
    id SERIAL PRIMARY KEY,
    counter INTEGER NOT NULL,
    "current_time" BIGINT NOT NULL,
    time TEXT NOT NULL,
    latitude REAL NOT NULL,
    longitude REAL NOT NULL,
    accuracy REAL,
    altitude REAL,
    traffic_total BIGINT,
    traffic_total_rx BIGINT,
    traffic_total_tx BIGINT
);

CREATE TABLE cells (
    id SERIAL PRIMARY KEY,
    measurement_id INTEGER NOT NULL,
    type TEXT,
    mcc INTEGER,
    mnc INTEGER,
    cell_identity BIGINT,
    tac INTEGER,
    earfcn INTEGER,
    band INTEGER,
    pci INTEGER,
    rsrp INTEGER,
    rsrq INTEGER,
    rssi INTEGER,
    rssnr INTEGER,
    asu_level INTEGER,
    cqi INTEGER,
    timing_advance INTEGER,
    FOREIGN KEY (measurement_id) REFERENCES based(id)
);

CREATE INDEX idx_based_time ON based(time);
CREATE INDEX idx_based_location ON based(latitude, longitude);
CREATE INDEX idx_cells_measurement ON cells(measurement_id);
CREATE INDEX idx_cells_mcc_mnc ON cells(mcc, mnc);
CREATE INDEX idx_cells_type ON cells(type);
\end{lstlisting}

\subsection{Реализация}

Сервер на C++ подключается к PostgreSQL через библиотеку libpq. В классе Database реализованы методы connect(), insertMeasurement(), insertCell().

\begin{lstlisting}[language=C++, caption=Класс Database для работы с PostgreSQL, label=lst:database_class]
class Database{
public:
    Database():conn(nullptr){}
    ~Database(){if (conn) PQfinish(conn);}
    bool connect(){
        const char* info="host=" DB_HOST " port=" DB_PORT 
                          " dbname=" DB_NAME " user=" DB_USER 
                          " password=" DB_PASSWORD;
        conn=PQconnectdb(info);
        if (PQstatus(conn)!=CONNECTION_OK){
            std::cerr<<"ОШИБКА подключения: "<<PQerrorMessage(conn)<<"\n";
            return false;}return true;}

    int insertMeasurement(const LocationData& data,int counter);
    bool insertCell(int measurement_id,const CellTowerData& cell);};
\end{lstlisting}

При получении JSON от Android сервер парсит данные, вызывает insertMeasurement для записи в таблицу based, а затем для каждой соты - insertCell.

\begin{lstlisting}[language=C++, caption=Вставка данных в БД, label=lst:db_insert]
int Database::insertMeasurement(const LocationData& data,int counter){
    std::ostringstream query;
    query << "INSERT INTO based (counter, current_time, time, "
          << "latitude, longitude, accuracy, altitude, "
          << "traffic_total, traffic_total_rx, traffic_total_tx) "
          << "VALUES ($1, $2, $3, $4, $5, $6, $7, $8, $9, $10) "
          << "RETURNING id";
    return measurement_id;}

bool Database::insertCell(int measurement_id,const CellTowerData& cell){
    std::ostringstream query;
    query << "INSERT INTO cells (measurement_id, type, mcc, mnc, "
          << "cell_identity, tac, earfcn, band, pci, rsrp, rsrq, "
          << "rssi, rssnr, asu_level, cqi, timing_advance) "
          << "VALUES ($1, $2, $3, $4, $5, $6, $7, $8, $9, $10, "
          << "$11, $12, $13, $14, $15, $16)";}
\end{lstlisting}

В GUI-потоке данные о сигналах накапливаются в структуре SignalHistory с мьютексом. Графики используют библиотеку ImPlot.

\begin{lstlisting}[language=C++, caption=Структура для хранения истории сигналов, label=lst:signal_history]
struct SignalHistory{
    std::vector<double> timestamps;
    std::map<int,std::vector<double>> rsrp_values;
    std::map<int,std::vector<double>> rsrq_values;
    std::map<int,std::vector<double>> rssi_values;
    std::map<int,std::vector<double>> rssnr_values;
    std::map<int,std::vector<double>> ss_rsrp_values;
    std::map<int,std::vector<double>> ss_rsrq_values;
    std::map<int,std::vector<double>> ss_sinr_values;
    std::map<int,std::vector<double>> dbm_values;
    std::mutex mutex;};
\end{lstlisting}

\begin{figure}[h]
\centering
\includegraphics[width=0.8\textwidth]{lab13_backend.png}
\caption{Окно сервера на C++ 2 версии}
\end{figure}

\section{Лабораторная работа №14. Отрисовка OpenStreetMap карты}

\subsection{Теория: проекция Меркатора и тайлы}

Для отображения карты на экране используется проекция Меркатора. Все карты OpenStreetMap разбиты на квадратные тайлы размером 256 × 256 пикселей. Каждый тайл однозначно идентифицируется уровнем масштаба zoom и координатами x, y.

\subsection{Реализация}

\subsubsection{Класс TileManager}

Создан класс TileManager, который управляет загрузкой, кэшированием и отображением тайлов OSM. Основные возможности: асинхронная загрузка тайлов из сети, кэширование на диске, декодирование PNG в RGBA с использованием libpng, создание текстур OpenGL.

\begin{lstlisting}[language=C++, caption=Метод загрузки тайла, label=lst:tile_download]
bool TileManager::downloadTile(int zoom,int x,int y,std::vector<uint8_t>& rgbaData,int& width,int& height){
    std::string url="https://tile.openstreetmap.org/"+std::to_string(zoom)+"/"+std::to_string(x)+"/"+std::to_string(y)+".png";
    CURL* curl=curl_easy_init();
    curl_easy_setopt(curl,CURLOPT_URL,url.c_str());
    curl_easy_setopt(curl,CURLOPT_WRITEFUNCTION,WriteCallback);
    CURLcode res=curl_easy_perform(curl);
    //обработка ответа и декодирование png
    return decodePngToRgba(response,rgbaData,width,height);}
\end{lstlisting}

\subsubsection{Пересчёт координат и подгрузка тайлов}

В GUI-потоке при отображении карты определяются текущие границы видимой области. На основе этих границ вычисляется оптимальный уровень масштаба zoom.

\begin{lstlisting}[language=C++, caption=Пересчёт географических координат в тайлы, label=lst:coord_to_tile]
int TileManager::getZoomForLimits(double minLon,double maxLon){
    double diff=maxLon-minLon;
    if (diff>90.0) return 4;
    else if (diff>45.0) return 5;
    else if (diff>22.5) return 6;
    else if (diff>11.25) return 7;
    else return 15;}

double TileManager::mercatorXToTileX(double mercatorX,int zoom){
    return (0.5+mercatorX/360.0)*(1<<zoom);}

double TileManager::mercatorYToTileY(double mercatorY,int zoom){
    double lat_rad=mercatorY*M_PI/180.0;
    double y_normalized=0.5-log(tan(M_PI/4.0+lat_rad/2.0))/(2.0*M_PI);
    return y_normalized*(1<<zoom);}
\end{lstlisting}

\subsubsection{Отображение карты в ImPlot}

Карта рисуется с помощью функции ImPlot::PlotImage(), которая принимает текстурный идентификатор и координаты углов тайла.

\begin{lstlisting}[language=C++, caption=Рендеринг тайлов карты, label=lst:tile_render]
void TileManager::renderTiles(double minX,double maxX,double minY,double maxY){
    int zoom=getZoomForLimits(minX,maxX);
    
    for (int x=minTileX;x<=maxTileX;x++){
        for (int y=minTileY;y<=maxTileY;y++){
            GLuint gpuId=getTextureId(zoom,x,y);
            if (gpuId!=0){
                ImPlotPoint minPoint{tileXToMercatorX(x,zoom),tileYToMercatorY(y+1,zoom)};
                ImPlotPoint maxPoint{tileXToMercatorX(x+1,zoom),tileYToMercatorY(y,zoom)};
                ImPlot::PlotImage(("##tile_"+tileId).c_str(),(ImTextureID)(intptr_t)gpuId,minPoint, maxPoint);}}}}
\end{lstlisting}

\subsubsection{Асинхронная загрузка}

Для плавной работы интерфейса загрузка тайлов выполняется в отдельном потоке через очередь задач.

\begin{lstlisting}[language=C++, caption=Асинхронная загрузка тайлов, label=lst:async_load]
void TileManager::fetchWorker(){
    while (m_running){
        TileJob job;
        bool hasJob=false;
        {   std::lock_guard<std::mutex> lock(m_jobMutex);
            if (!m_jobQueue.empty()){
                job = m_jobQueue.front();
                m_jobQueue.pop();
                hasJob=true;}}
        if (hasJob){
            std::vector<uint8_t> rgbaData;
            int width, height;
            if (downloadTile(job.zoom,job.x,job.y,rgbaData,width,height)){}
        }else{std::this_thread::sleep_for(std::chrono::milliseconds(10));}}}
\end{lstlisting}

\begin{figure}[h]
\centering
\includegraphics[width=0.8\textwidth]{lab14_backend.png}
\caption{Окно сервера на C++ 3 версии}
\end{figure}

\section{Лабораторная работа №15. Генерация тепловой карты}

\subsection{Теория: метод обратных взвешенных расстояний (IDW)}

Тепловая карта строится на основе интерполяции экспериментальных измерений. Используется метод обратных взвешенных расстояний (IDW, Inverse Distance Weighting):

\[
\hat{Z}(x_0, y_0) = \frac{\sum_{i=1}^{n} \frac{Z_i}{d_i^p}}{\sum_{i=1}^{n} \frac{1}{d_i^p}}
\]

где \(d_i\) — расстояние до i-й точки измерения, \(p\) — параметр степени (в проекте \(p=2\)).

Критерии оценки RSRP:
\begin{itemize}
    \item Excellent: > -80 dBm (красный цвет)
    \item Good: от -80 до -90 dBm
    \item Fair/Marginal: от -90 до -100 dBm
    \item Poor/Weak: от -100 до -110 dBm (тёмно-синий)
    \item No Signal: < -110 dBm (не закрашивается)
\end{itemize}

\subsection{Реализация}

\subsubsection{Класс HeatmapData}

Создан класс HeatmapData, который загружает из базы данных PostgreSQL все измерения LTE. Данные группируются по метрике и по EARFCN. Для каждого набора вычисляется динамический диапазон значений.

\begin{lstlisting}[language=C++, caption=Загрузка данных для тепловой карты из БД, label=lst:heatmap_load]
std::vector<HeatmapDbRow> Database::getHeatmapRows(){
    std::vector<HeatmapDbRow> rows;
    const char* query=
        "SELECT b.latitude, b.longitude, b.accuracy, b.altitude, "
        "c.earfcn, c.rsrp, c.rsrq, c.rssi "
        "FROM based b "
        "JOIN cells c ON c.measurement_id = b.id "
        "WHERE c.type = 'LTE' "
        "ORDER BY b.id";
    
    PGresult* res=PQexec(conn,query);
    for (int i=0;i<PQntuples(res);++i){
        HeatmapDbRow row;
        row.latitude=std::stod(PQgetvalue(res,i,0));
        row.longitude=std::stod(PQgetvalue(res,i,1));
        row.earfcn=std::stoi(PQgetvalue(res,i,4));
        row.rsrp=std::stoi(PQgetvalue(res,i,5));
        rows.push_back(row);}return rows;}
\end{lstlisting}

\subsubsection{Рендеринг тайлов тепловой карты}

Метод renderTile() для заданного тайла создаёт изображение 256×256 пикселей. Для каждого пикселя вычисляются географические координаты, затем по методу IDW интерполируется значение.

\begin{lstlisting}[language=C++, caption=Интерполяция IDW для тепловой карты, label=lst:idw_interpolation]
bool HeatmapData::renderTile(const HeatmapSelection& selection,int zoom,int x,int y,std::vector<uint8_t>& rgbaData,int& width,int& height) const{
    width=256;height=256;
    rgbaData.assign(width*height*4,0);
    
    for (int pixelY=0;pixelY<height;++pixelY){
        const double latitude=tilePixelToLatitude(zoom,y,pixelY,height);
        for (int pixelX=0;pixelX<width;++pixelX){
            const double longitude=tilePixelToLongitude(zoom,x,pixelX,width);
            double weightedSum=0.0;
            double totalWeight=0.0;
            for (const Sample* sample:candidates){
                const double distance=haversineMeters(latitude,longitude,sample->latitude,sample->longitude);
                const double weight=1.0/std::pow(distance,selection.power);
                weightedSum+=sample->value*weight;
                totalWeight+=weight;}
            
            if (totalWeight>0.0){
                double interpolatedValue=weightedSum/totalWeight;}}}
    return hasVisiblePixels;}
\end{lstlisting}

\subsubsection{Цветовая схема}

Для визуализации значений используется градиент от синего (худшие значения) через зелёный и жёлтый к красному (лучшие значения).

\begin{lstlisting}[language=C++, caption=Функция цветовой схемы, label=lst:colorize]
void HeatmapData::colorize(double ratio,float opacity,uint8_t& r,uint8_t& g,uint8_t& b,uint8_t& a){
    static const ColorStop stops[]={
        {0.0,10.0,20.0,120.0},
        {0.35,0.0,120.0,255.0},
        {0.7;255.0,200.0,0.0},
        {1.0,255.0,0.0,0.0}};
    r=static_cast<uint8_t>(left->r+(right->r-left->r)*localRatio);
    g=static_cast<uint8_t>(left->g+(right->g-left->g)*localRatio);
    b=static_cast<uint8_t>(left->b+(right->b-left->b)*localRatio);
    a=static_cast<uint8_t>(255.0f*opacity);}
\end{lstlisting}

\subsubsection{Интеграция в GUI}

В окне ImGui добавлена панель управления тепловой картой с выбором метрики, EARFCN и параметров интерполяции.

\begin{lstlisting}[language=C++, caption=Панель управления тепловой картой в GUI, label=lst:heatmap_gui]
ImGui::Checkbox("Show heatmap",&heatmapEnabled);
const char* metricItems[]={"RSRP","RSRQ","RSSI","Altitude"};
ImGui::Combo("Criterion",&selectedMetric,metricItems,IM_ARRAYSIZE(metricItems));

std::vector<int> availableEarfcns=g_tileManager.getAvailableHeatmapEarfcns(metric);
if (ImGui::BeginCombo("EARFCN",std::to_string(selectedEarfcn).c_str())){
    for (int earfcn:availableEarfcns){
        if (ImGui::Selectable(std::to_string(earfcn).c_str(),selectedEarfcn==earfcn)){
            selectedEarfcn=earfcn;
            heatmapGenerated=false;}}ImGui::EndCombo();}

if (ImGui::Button("Generate heatmap")){
    heatmapGenerated=true;
    g_tileManager.clearHeatmapCache();}

g_tileManager.setHeatmapSelection({
    heatmapEnabled && heatmapGenerated,
    metric, selectedEarfcn,radiusMeters,
    displayRadiusMeters,2.0,0.55f});
\end{lstlisting}

\begin{figure}[h]
\centering
\includegraphics[width=0.8\textwidth]{lab15_backend.png}
\caption{Окно сервера на C++ 4 версии}
\end{figure}

\endinput